```latex
\[
\newcommand{\R}{\mathbb{R}}
\newcommand{\diag}{\operatorname{diag}}
\]

Let
\[
\|A\|_{r\to r}:=\sup_{x\neq 0}\frac{\|Ax\|_r}{\|x\|_r}
\]
denote the induced \(\ell^r\)-operator norm, and let
\[
\rho_r(A):=\inf\left\{\|DAD^{-1}\|_{r\to r}:D\ \text{is positive diagonal}\right\}.
\]
Here a positive diagonal matrix means an invertible diagonal matrix whose diagonal entries are all positive.

We prove the two claims.

\[
\textbf{(a) Proof that }\rho_r(M_n)=\rho_r(H).
\]

First note the following elementary fact. If \(A\in \R^{m\times m}\) and \(s\geq 0\), then
\[
\left\|\diag(A,0_s)\right\|_{r\to r}=\|A\|_{r\to r}.
\]
Indeed, for \(x=(x_1,x_2)\in \R^m\oplus \R^s\),
\[
\left\|\diag(A,0_s)x\right\|_r
=
\|(Ax_1,0)\|_r
=
\|Ax_1\|_r
\leq
\|A\|_{r\to r}\|x_1\|_r
\leq
\|A\|_{r\to r}\|x\|_r.
\]
Thus
\[
\left\|\diag(A,0_s)\right\|_{r\to r}\leq \|A\|_{r\to r}.
\]
Conversely, since the unit sphere in \(\R^m\) is compact and \(x_1\mapsto \|Ax_1\|_r\) is continuous, there exists \(x_1^\ast\in \R^m\) with \(\|x_1^\ast\|_r=1\) such that
\[
\|Ax_1^\ast\|_r=\|A\|_{r\to r}.
\]
Testing \(\diag(A,0_s)\) on \(x=(x_1^\ast,0)\) gives
\[
\left\|\diag(A,0_s)\right\|_{r\to r}
\geq
\left\|\diag(A,0_s)(x_1^\ast,0)\right\|_r
=
\|Ax_1^\ast\|_r
=
\|A\|_{r\to r}.
\]
Therefore
\[
\left\|\diag(A,0_s)\right\|_{r\to r}=\|A\|_{r\to r}.
\]

Now let \(D\in \R^{n\times n}\) be positive diagonal. With respect to the splitting
\[
\R^n=\R^m\oplus \R^{n-m},
\]
we can write
\[
D=\diag(D_1,D_2),
\]
where \(D_1\in \R^{m\times m}\) and \(D_2\in \R^{(n-m)\times(n-m)}\) are positive diagonal. Then
\[
D M_n D^{-1}
=
\diag(D_1,D_2)
\begin{pmatrix}
H&0\\
0&0_{n-m}
\end{pmatrix}
\diag(D_1^{-1},D_2^{-1})
=
\diag(D_1HD_1^{-1},0_{n-m}).
\]
By the norm identity just proved,
\[
\|D M_n D^{-1}\|_{r\to r}
=
\|D_1HD_1^{-1}\|_{r\to r}.
\]
Taking the infimum over all positive diagonal \(D=\diag(D_1,D_2)\) therefore gives
\[
\rho_r(M_n)
=
\inf_{D_1,D_2}\|D_1HD_1^{-1}\|_{r\to r}
=
\inf_{D_1}\|D_1HD_1^{-1}\|_{r\to r}
=
\rho_r(H).
\]
Thus
\[
\boxed{\rho_r(M_n)=\rho_r(H)}
\]
for every \(r\geq 1\).

\[
\textbf{(b) Proof of Conjecture }3.
\]

We use the following standard fact for Hadamard matrices. If \(H\in \R^{m\times m}\) is a Hadamard matrix, so that
\[
HH^{\mathsf T}=mI_m,
\]
then for every \(q\geq 2\),
\[
\rho_q(H)=m^{1-\frac1q}.
\]

For completeness, we recall the proof. First, since \(D=I_m\) is an admissible positive diagonal matrix,
\[
\rho_q(H)\leq \|H\|_{q\to q}.
\]
Let \(x\in \R^m\) and set \(y=Hx\). Since \(H^{\mathsf T}H=mI_m\),
\[
\|y\|_2=\|Hx\|_2=\sqrt m\,\|x\|_2
\leq
\sqrt m\,m^{\frac12-\frac1q}\|x\|_q
=
m^{1-\frac1q}\|x\|_q.
\]
Also, because every entry of \(H\) is \(\pm 1\), each coordinate of \(Hx\) satisfies
\[
|(Hx)_i|\leq \sum_{j=1}^m |x_j|=\|x\|_1
\leq
m^{1-\frac1q}\|x\|_q.
\]
Thus
\[
\|Hx\|_\infty\leq m^{1-\frac1q}\|x\|_q.
\]
For \(q\geq 2\),
\[
\|y\|_q^q
=
\sum_{i=1}^m |y_i|^q
\leq
\|y\|_\infty^{q-2}\|y\|_2^2.
\]
Combining the previous two bounds gives
\[
\|Hx\|_q
\leq
m^{1-\frac1q}\|x\|_q.
\]
Therefore
\[
\|H\|_{q\to q}\leq m^{1-\frac1q},
\]
and hence
\[
\rho_q(H)\leq m^{1-\frac1q}.
\]

For the reverse inequality, let \(D=\diag(d_1,\dots,d_m)\) be any positive diagonal matrix. Choose an index \(i\) such that
\[
d_i=\max_{1\leq j\leq m} d_j.
\]
Let \(u^{(i)}\in \{\pm 1\}^m\) be the transpose of the \(i\)-th row of \(H\). Since the rows of \(H\) are mutually orthogonal,
\[
H u^{(i)} = m e_i.
\]
Set
\[
x=D u^{(i)}.
\]
Then
\[
DHD^{-1}x
=
DHu^{(i)}
=
D(me_i)
=
m d_i e_i.
\]
Moreover,
\[
\|x\|_q
=
\|D u^{(i)}\|_q
=
\left(\sum_{j=1}^m d_j^q\right)^{1/q}
\leq
\left(m d_i^q\right)^{1/q}
=
m^{1/q}d_i.
\]
Hence
\[
\|DHD^{-1}\|_{q\to q}
\geq
\frac{\|DHD^{-1}x\|_q}{\|x\|_q}
\geq
\frac{m d_i}{m^{1/q}d_i}
=
m^{1-\frac1q}.
\]
Since this holds for every positive diagonal \(D\),
\[
\rho_q(H)\geq m^{1-\frac1q}.
\]
Therefore
\[
\boxed{\rho_q(H)=m^{1-\frac1q}}
\]
for every \(q\geq 2\).

Hadamard matrices of orders \(4\) and \(8\) exist by the Sylvester construction, which recursively defines
\[
H_{2t}
=
\begin{pmatrix}
H_t&H_t\\
H_t&-H_t
\end{pmatrix}.
\]
This standard construction yields Hadamard matrices in every order \(2^k\). ([en.wikipedia.org](https://en.wikipedia.org/wiki/Hadamard_matrix))

Now let \(n>6\) and let \(p>2\) be an integer.

\[
\textbf{Case }1:\ n=7.
\]
Take
\[
m=4
\]
and choose the \(4\times 4\) Hadamard matrix
\[
H_4=
\begin{pmatrix}
1&1&1&1\\
1&-1&1&-1\\
1&1&-1&-1\\
1&-1&-1&1
\end{pmatrix}.
\]
Define
\[
M_7=\diag(H_4,0_3)\in \R^{7\times 7}.
\]
By part \(\text{(a)}\),
\[
\rho_p(M_7)=\rho_p(H_4)
\qquad\text{and}\qquad
\rho_{p+1}(M_7)=\rho_{p+1}(H_4).
\]
Using the Hadamard formula with \(m=4\),
\[
\rho_p(M_7)=4^{1-\frac1p},
\]
and
\[
\rho_{p+1}(M_7)=4^{1-\frac1{p+1}}.
\]
Since
\[
\frac1p>\frac1{p+1},
\]
we have
\[
1-\frac1p<1-\frac1{p+1}.
\]
Because \(4>1\), it follows that
\[
4^{1-\frac1p}<4^{1-\frac1{p+1}}.
\]
Therefore
\[
\boxed{\rho_p(M_7)<\rho_{p+1}(M_7)}.
\]

\[
\textbf{Case }2:\ n\geq 8.
\]
Take
\[
m=8
\]
and let \(H_8\) be an \(8\times 8\) Hadamard matrix. Define
\[
M_n=\diag(H_8,0_{n-8})\in \R^{n\times n}.
\]
By part \(\text{(a)}\),
\[
\rho_p(M_n)=\rho_p(H_8)
\qquad\text{and}\qquad
\rho_{p+1}(M_n)=\rho_{p+1}(H_8).
\]
Using the Hadamard formula with \(m=8\),
\[
\rho_p(M_n)=8^{1-\frac1p},
\]
and
\[
\rho_{p+1}(M_n)=8^{1-\frac1{p+1}}.
\]
Again,
\[
1-\frac1p<1-\frac1{p+1},
\]
and since \(8>1\),
\[
8^{1-\frac1p}<8^{1-\frac1{p+1}}.
\]
Thus
\[
\boxed{\rho_p(M_n)<\rho_{p+1}(M_n)}.
\]

Combining the two cases, for every integer \(n>6\) and every integer \(p>2\), there exists a matrix \(M\in \R^{n\times n}\) such that
\[
\rho_p(M)<\rho_{p+1}(M).
\]
This proves Conjecture \(3\).
```